% Appendix A -- Messages (PUB-REF-0001 AN-03; IBM 3.10 Reference Appendix A,
% ibm-dos.md Sec 8.5).  The message lists are GENERATED from the locked
% catalogue spec/dos_messages.json by tools/pubs_msgtable.c (device / other
% selectors); only the Explanation and Action of each message are authored
% here, keyed by catalogue id with \MessageDetail; a missing or stray detail
% fails the build (CP-06).  Where each message comes from (Law 1):
% MSG-DOS-0001 int21.c int24_dispatch_body (A/R/F, crit_error_action);
% 0002..0005, 0007, 0009..0012, 0016..0020 command.c (and the utilities);
% 0006, 0008, 0013, 0014, 0015 are catalogue entries no code path of this
% version displays (DRAFT, each says so).  0027..0037 command.c (ECHO,
% BREAK, SET/PROMPT/PATH and EXEC environment, DATE/TIME, FOR/IF/CALL
% nesting BATCH_NEST_MAX 8, GOTO); 0035 also edlin.c CAP; 0038..0041
% dosutil.c (ATTRIB, LABEL) and xcopy_tree.c (TREE, XCOPY); 0042..0049
% chkdsk.c (0044 + path + 0045 is one message, joined by pubs_msgtable;
% incomplete tree => no lost-chain recovery); 0050, 0051 xcopy_tree.c;
% 0052..0058 edlin.c; 0059 find_program.asm (quoted string required); 0060
% sort_program.asm (INBUF_MAX 64 KiB); 0062 command.c builtin_ren
% (WL-0154; 0061 is the COPY footer fragment).  Rows 0027..0060 from WL-0142; the
% register is the IBM DOS 3.10 Reference Appendix A (pp.A-1..A-60), the
% wording is ours.  "Other Messages Displayed by Commands" quotes the
% strings of command.c and chkdsk.c that are not in the catalogue.
% GATES: test-dosmsg test-shell test-command test-attrib test-label-vol test-chkdsk test-xcopy-tree test-edlin-emu
\Gates{test-dosmsg, test-shell, test-command, test-attrib, test-label-vol, test-chkdsk, test-xcopy-tree, test-edlin-emu}
\chapter{Messages}
\MessageDetailsOn

\section{How to Use This Appendix}
This appendix lists the messages InitechDOS displays.  Each message is shown
exactly as it appears on the screen, followed by an explanation of what it
means and the action you should take.  Where a message names a drive, the
letter \ph{d} stands for the drive letter.  The number beside each message
identifies it when you call the Help Desk (extension 2504).

The device error messages come first.  The other messages follow in
alphabetical order, and then the messages that individual commands display.

\section{Responses}
When InitechDOS cannot read or write a disk, it displays a device error
message and asks \texttt{Abort, Retry, Fail?}  Press one of these letters;
you need not press \key{Enter}:

\SideHead{A (Abort)}Ends the program that was using the disk and returns to
the command prompt.

\SideHead{R (Retry)}Tries the operation again.  Correct the cause first,
for example by closing the drive door.

\SideHead{F (Fail)}Tells the program that the operation failed and lets it
continue.  The program decides what to do next.

Any other key displays the question again.

\section{Device Error Messages}
\MessageDetail{MSG-DOS-0001}{InitechDOS could not read or write a disk and
asks what to do.  It is displayed after a device error message.}{Press
\type{R} to try again after correcting the problem, \type{A} to end the
program, or \type{F} to let the program continue.}
\MessageDetail{MSG-DOS-0007}{InitechDOS could not read the disk in drive
\ph{d}.  \cmd{DIR} displays this message when it cannot read the volume
label of the disk.}{Make sure the disk is formatted and seated in the drive,
then try the command again.  If the message persists, the disk may be
damaged; copy the files you can read to another disk.}
\MessageDetail{MSG-DOS-0006}{The drive \ph{d} was not ready: the door may be
open or the drive may be empty.\draft{see bead initech-vlhp: catalogue entry not displayed by any
code path of this version}}{Put a disk in the drive, close the door, and
press \type{R}.}
\MessageDetail{MSG-DOS-0008}{InitechDOS tried to write on a disk in drive
\ph{d} that is write-protected.\draft{see bead initech-vlhp: catalogue entry not displayed by any
code path of this version}}{Remove the write protection, or use another
disk, and press \type{R}.}
\input{dos_messages_device}

\section{Other Messages}
\MessageDetail{MSG-DOS-0009}{The file is read-only, or the disk cannot be
written or read.  \cmd{DEL} and \cmd{COPY} display it for a read-only file;
\cmd{VOL} displays it when the disk cannot be read.}{Use \cmd{ATTRIB} to
display the file's attributes and, if you must change the file, \cmd{ATTRIB
-R} to make it read-write.}
\MessageDetail{MSG-DOS-0012}{\cmd{DEL} was asked to delete every file in a
directory (\file{*.*}).}{Type \type{Y} and press \key{Enter} to delete the
files, or any other answer to keep them.}
\MessageDetail{MSG-DOS-0002}{The command name is not an internal command,
and no program or batch file of that name was found in the current directory
or in the directories named by \texttt{PATH}.  \cmd{REN} and \cmd{COPY} also
display it when a file cannot be renamed or created.}{Check the spelling of
the command.  If it is a program, make sure its directory is the current
directory, is named in \texttt{PATH}, or is typed before the name.}
\MessageDetail{MSG-DOS-0014}{The command processor could not be found or
started.\draft{see bead initech-vlhp: catalogue entry not displayed by any code path of this
version: COMMAND.COM is part of the system}}{Restart the system.  If the
message appears again, call the Help Desk.}
\MessageDetail{MSG-DOS-0020}{The source and target named in \cmd{COPY} or
\cmd{XCOPY} are the same file.  Nothing was copied.}{Give the copy a
different name, or copy it to another directory.}
\MessageDetail{MSG-DOS-0003}{No file matches the name you gave.}{Check the
name and the directory.  Use \cmd{DIR} to list the files.}
\MessageDetail{MSG-DOS-0013}{The command processor must be loaded again from
the disk in drive \ph{d}.\draft{see bead initech-vlhp: catalogue entry not displayed by any code
path of this version: COMMAND.COM is never reloaded}}{Put the InitechDOS
disk in the drive and press any key.}
\MessageDetail{MSG-DOS-0005}{There is not enough room on the disk.  When you
use a pipe (\texttt{|}), InitechDOS displays this message if it cannot
create the temporary file that joins the two commands.}{Delete files you no
longer need, then try again.}
\MessageDetail{MSG-DOS-0018}{\cmd{CD} could not find the directory you
named.}{Check the path.  Use \cmd{DIR} or \cmd{TREE} to see the
directories on the disk.}
\MessageDetail{MSG-DOS-0004}{The drive letter is not a drive of this system.
In this version only drive~A exists.}{Use drive~A, or leave the drive
out.}
\MessageDetail{MSG-DOS-0019}{\cmd{RD} could not remove the directory: it does
not exist, it still holds files or subdirectories, or it is the current
directory.}{Delete the files in the directory, move to another directory,
and try again.}
\MessageDetail{MSG-DOS-0015}{The disk in drive~A cannot start the
system.\draft{see bead initech-vlhp: catalogue entry not displayed by any code path of this
version}}{Put the InitechDOS disk in the drive and restart the system.}
\MessageDetail{MSG-DOS-0016}{\cmd{DIR /P} has filled the screen.}{Press any
key to see the next screen of the listing.}
\MessageDetail{MSG-DOS-0011}{The command needs a parameter, such as a file
name, that you did not give.}{Type the command again with the missing
parameter.  The Format line of the command shows what is needed.}
\MessageDetail{MSG-DOS-0010}{The command was given more parameters than it
uses, or (\cmd{DIR}, \cmd{VOL}) a switch it does not have.}{Type the command again with only the parameters shown in its Format
line.}
\MessageDetail{MSG-DOS-0017}{\cmd{MD} could not create the directory: a file
or directory of that name already exists, a directory in the path does not
exist, or the disk is full.}{Choose another name, or check the path.}
% The catalogue rows MSG-DOS-0027..0060 (ADR-0003 Appendix C, the batch-2
% merge) are the strings command.c, dosutil.c, chkdsk.c, xcopy_tree.c,
% edlin.c, find_program.asm and sort_program.asm display; each names the
% command that displays it.  Three "Insufficient memory" rows and two
% "Invalid parameter(s)" / "Invalid number of parameters" rows are distinct
% catalogue ids for the same text, one per program, and are explained once
% each under their own number.  MSG-DOS-0045 is the tail of MSG-DOS-0044
% and is printed with it (tools/pubs_msgtable.c), explained under 0044.
\MessageDetail{MSG-DOS-0027}{\cmd{ECHO} was typed without a parameter.  It
displays whether the commands of a batch file are displayed as they are
carried out.  At the command prompt \cmd{ECHO} is always on.}{None.  To
change the setting in a batch file, use \cmd{ECHO OFF}.}
\MessageDetail{MSG-DOS-0028}{\cmd{ECHO} was typed without a parameter in a
batch file after \cmd{ECHO OFF}.  The commands of the batch file are not
displayed as they are carried out.}{None.  To display them again, use
\cmd{ECHO ON}.}
\MessageDetail{MSG-DOS-0029}{The CTRL-BREAK check is on.  \cmd{BREAK}
displays the setting when it is typed with no parameter, and again after
\kw{ON} or \kw{OFF}.}{None.  To turn the check off, use \cmd{BREAK OFF}.}
\MessageDetail{MSG-DOS-0030}{The CTRL-BREAK check is off.  \cmd{BREAK}
displays the setting when it is typed with no parameter, and again after
\kw{ON} or \kw{OFF}.}{None.  To turn the check on, use \cmd{BREAK ON}.}
\MessageDetail{MSG-DOS-0031}{The environment is full.  \cmd{SET},
\cmd{PROMPT} and \cmd{PATH} display this message when the new setting does
not fit; the setting is not changed.  It is also displayed when a program
is started and its copy of the environment cannot be made; the program is
given an empty environment.}{Remove variables you no longer need with
\cmd{SET} \ph{name}\texttt{=}, then try again.}
\MessageDetail{MSG-DOS-0032}{\cmd{SET} was asked to display a variable that
is not in the environment.}{Check the name.  Use \cmd{SET} alone to list
the variables.}
\MessageDetail{MSG-DOS-0033}{The date typed for \cmd{DATE}, or at the
start-up prompt, is not a real date or is not in the form
\ph{mm}\texttt{-}\ph{dd}\texttt{-}\ph{yy}.  The date is not changed.}{Type
the date again in the form shown.}
\MessageDetail{MSG-DOS-0034}{The time typed for \cmd{TIME}, or at the
start-up prompt, is not in the form \ph{hh}\texttt{:}\ph{mm}, or a number
is out of range.  The time is not changed.}{Type the time again in the form
shown.}
\MessageDetail{MSG-DOS-0035}{\cmd{EDLIN}: the lines do not fit in memory;
the command is not carried out.  \cmd{FOR}: the file names that match the
set do not fit in the space \cmd{FOR} keeps for them.}{\cmd{EDLIN}: write
part of the file to disk with \kw{W}, then go on.  \cmd{FOR}: name fewer
files in the set, or use more than one \cmd{FOR}.}
\MessageDetail{MSG-DOS-0036}{The batch file has no label of the name given
in \cmd{GOTO}.  The batch file ends.}{Correct the label or the \cmd{GOTO}
line.}
\MessageDetail{MSG-DOS-0037}{Batch files have called each other with
\cmd{CALL}, or \cmd{IF} and \cmd{FOR} commands have been put inside each
other, more than eight levels deep.  The command is not carried
out.}{Reduce the number of levels.}
\MessageDetail{MSG-DOS-0038}{\cmd{ATTRIB}, \cmd{TREE} or \cmd{XCOPY} was
given too many or too few file names.}{Check the Format line of the
command.}
\MessageDetail{MSG-DOS-0039}{\cmd{ATTRIB}, \cmd{TREE} or \cmd{XCOPY} could
not find the file or directory, or a directory was named where a file was
expected.}{Check the name and the path.}
\MessageDetail{MSG-DOS-0040}{\cmd{LABEL} found no free entry in the root
directory for the volume label.}{Delete a file from the root directory and
try again.}
\MessageDetail{MSG-DOS-0041}{The label typed for \cmd{LABEL} contains a
character that is not allowed in a volume label, or has no
characters.}{Type a label of letters, digits and spaces.}
\MessageDetail{MSG-DOS-0042}{\cmd{CHKDSK} found that the file named before
this message contains a cluster number that is not valid.}{Run
\cmd{CHKDSK /F} to end the file at the last valid cluster.  Without
\kw{/F}, nothing is changed.}
\MessageDetail{MSG-DOS-0043}{\cmd{CHKDSK} found the cluster \ph{n}
allocated to two files.  The message is displayed after the name of each of
the two files.}{\cmd{CHKDSK} does not correct this, even with \kw{/F}.
Copy both files, delete the originals, and check the copies.}
\MessageDetail{MSG-DOS-0044}{\cmd{CHKDSK} could not move into the
subdirectory \ph{path}, because its directory entry names a cluster that is
not valid.  The files below it are not checked, and lost clusters are not
recovered.}{Copy the files you can read to another disk.}
\MessageDetail{MSG-DOS-0046}{\cmd{CHKDSK /F} could not create a
\file{FILE}\ph{nnnn}\file{.CHK} file for a lost chain, usually because the
root directory is full.}{Delete or move files from the root directory and
run \cmd{CHKDSK /F} again.}
\MessageDetail{MSG-DOS-0047}{\cmd{CHKDSK} found errors, but \kw{/F} was not
given, so nothing was changed.}{Run \cmd{CHKDSK /F} to correct the errors.}
\MessageDetail{MSG-DOS-0048}{\cmd{CHKDSK} was given a switch other than
\kw{/F} or \kw{/V}, or more than one file name.}{Check the Format line of
\cmd{CHKDSK}.}
\MessageDetail{MSG-DOS-0049}{\cmd{CHKDSK} could not read or write the file
allocation table, or the disk is not formatted as InitechDOS
expects.}{Copy the files you can read to another disk.}
\MessageDetail{MSG-DOS-0050}{\cmd{XCOPY} was asked to copy a directory, with
its subdirectories, into itself or into one of its own
subdirectories.}{Copy the directory to a place outside it.}
\MessageDetail{MSG-DOS-0051}{\cmd{XCOPY} or \cmd{TREE} was given a switch it
does not accept, or switches that cannot be used together.}{Check the
Format line of the command.}
\MessageDetail{MSG-DOS-0052}{The line is not a valid \cmd{EDLIN} command:
an unknown letter, a line number out of range, or lines that overlap the
destination of \kw{C} or \kw{M}.  It is also displayed when \cmd{EDLIN} is
started without a file name.}{Type the command again.}
\MessageDetail{MSG-DOS-0053}{\cmd{EDLIN} could not read, create or rename a
file, for example because it is read-only.}{Check the file's attributes
with \cmd{ATTRIB}.}
\MessageDetail{MSG-DOS-0054}{The disk filled while \cmd{EDLIN} was writing
the file.  The original file is unchanged.}{Make room on the disk, or quit
with \kw{Q}.}
\MessageDetail{MSG-DOS-0055}{A replacement made with \kw{R} would make the
line longer than \cmd{EDLIN} allows.  The line is not changed.}{Replace a
shorter text, or split the line.}
\MessageDetail{MSG-DOS-0056}{\kw{S} or \kw{R} found no match in the lines
searched.}{Check the text; capital and small letters must match.}
\MessageDetail{MSG-DOS-0057}{\cmd{EDLIN} does not edit a file with the
extension \file{.BAK}; that name is kept for the backup copy it
makes.}{Rename the file with \cmd{REN} and edit it under the new name.}
\MessageDetail{MSG-DOS-0058}{The file's \file{.BAK} file is read-only, so
\kw{E} will not replace it; the original is not kept.}{Make the
\file{.BAK} file read-write with \cmd{ATTRIB -R} first if you want a
backup.}
\MessageDetail{MSG-DOS-0059}{\cmd{FIND} was not given the string to search
for.  The string must be enclosed in quotation marks.}{Type the command
again with the string in quotation marks.}
\MessageDetail{MSG-DOS-0060}{The input to \cmd{SORT} is larger than the
64K bytes it can sort.  Nothing is written.}{Sort a smaller file, or
divide the file into parts.}
\MessageDetail{MSG-DOS-0062}{\cmd{REN} could not rename a file: a file
with the new name already exists, or no file matches the name you gave.
When a wildcard names several files, the others are still renamed.}{Choose
another new name, or use \cmd{DIR} to check the name of the file to
rename.}
\input{dos_messages_other}
\CheckMessageDetails

\section{Other Messages Displayed by Commands}
These messages are displayed by particular commands and are not numbered.
The command is shown at the right.

\OtherMessage{All specified file(s) are contiguous}{CHKDSK}{Every file you
named is stored in one piece.}{None.}
\OtherMessage{No Path}{PATH}{No search path is set: only the current
directory is searched for commands.}{Set one with \cmd{PATH} if you need
it.}
\OtherMessage{Strike a key when ready . . .}{PAUSE}{A batch file is paused at
a \cmd{PAUSE} subcommand.}{Do what the batch file asks, then press any key.
Press \key{Ctrl-C} to end the batch file.}
